\documentclass[10pt,a4paper]{amsart}
\usepackage[margin=19mm]{geometry}
\usepackage{amsmath,amssymb,mathtools}
\usepackage[scr=rsfs,cal=euler]{mathalfa}
\usepackage{xfrac}
\usepackage[unicode,hidelinks,bookmarks=false]{hyperref}
\newcommand{\ie}{i.e.\ }
\newcommand{\cf}{cf.\ }
\newcommand{\resp}{respectively\ }
\newcommand{\Subsection}{\subsection}
\providecommand{\nobreakdash}{\nobreak\hbox{-}}
\setlength{\emergencystretch}{2em}
\multlinegap0pt
\usepackage{dsfont}
\usepackage{bm}
\usepackage{cancel}
\usepackage{amssymb}
\usepackage{mathtools}
\usepackage{xcolor}
\newtheorem{lemma}{Lemma}
\newcommand{\E}{\mathbb{E}}
\newcommand{\Ll}{\left}
\newcommand{\N}{\mathbb{N}}
\newcommand{\R}{\mathbb{R}}
\newcommand{\Rr}{\right}
\renewcommand{\bar}{\overline}
\newcommand{\cM}{\mathcal{M}}
\renewcommand{\d}{\mathrm{d}}
\renewcommand{\geq}{\geqslant}
\renewcommand{\leq}{\leqslant}
\newcommand{\prog}{\mathsf{Prog}}
\DeclareMathOperator{\sign}{sign}
\title{One-sided large deviations for the ground-state energy of spin glasses}
\author{Hong-Bin Chen, Alice Guionnet, Justin Ko, Bertrand Lacroix-A-Chez-Toine, Jean-Christophe Mourrat}
\date{Source: arXiv:2603.06368v1 (CC BY 4.0)}
\begin{document}
\maketitle

\noindent\emph{Source and licence.} One-sided large deviations for the ground-state energy of spin glasses. Version arXiv:2603.06368v1 (CC BY 4.0), distributed by arXiv under the Creative Commons Attribution 4.0 International licence, http://creativecommons.org/licenses/by/4.0/, as recorded in the arXiv licence metadata for this exact version and shown on the abstract page https://arxiv.org/abs/2603.06368v1. Full licence notice: \texttt{licenses/arxiv-2603.06368-CC-BY-4.0.txt}.\par\smallskip

\noindent\textit{Editorial scope.} Counted unit: the article's lemma l.laplace\_> (source0.tex lines 892--898). The article's complete original proof of it is delivered with it (source0.tex lines 901--977, one \texttt{proof} environment of 77 lines). No mathematics is written, repaired or re-derived: the statement and the proof are the article's own lines, in the article's own notation, and no retained line is substituted or re-flowed.  Retained uncounted as the source-local context the proof needs: lines 689--691; lines 707--709; lines 854--856; lines 859--861; lines 867--869; lines 874--878; lines 880--886.  Nothing was omitted from any retained span, so the package carries no omission record.  No retained line contains a citation, so the wrapper carries no bibliography and no bibliography entry.  No retained line contains a figure environment, an image-inclusion command or drawing code, so no image is delivered and the package declares no asset dependency.  Editorial formatting only, in addition: an AMS \texttt{amsart} preamble that restates the article's own declarations and macros used by the retained lines, namely the environments \texttt{lemma} on separate counters, together with the standard packages those lines need.  The retained environments print in this document as Lemma 1 in order of appearance, whereas the article prints the same items with the numbering of its own full document.  The complete native source of the article as distributed in the arXiv e-print for this exact version is bundled unchanged as \texttt{sources/195861/source0.tex}.
\begin{align}\label{e.lim1/sNloge^sNL_N}
    \lim_{N\to\infty}\frac{1}{sN}\log \E\Ll[e^{sNL_N}\Rr] = \lim_{\beta\to \infty}\lim_{N\to\infty}\frac{1}{sN}\log \E\Ll[Z_N(\beta)^\frac{s}{\beta}\Rr].
\end{align}

\begin{align}\label{e.lim1/sNlogE[Z_N(beta)^s/beta]}
    \lim_{N\to\infty}\frac{1}{sN}\log \E\Ll[Z_N(\beta)^\frac{s}{\beta}\Rr] = \inf_{\gamma \in \cM_s^\beta}\bigg( \frac{1}{s}\log\E\Ll[e^{s\Psi_{\gamma,\beta}(0,g_1+h_1)}\Rr] - \frac{1}{2} \int_0^1 t \xi''(t) \gamma(t) \, \d t \bigg).
\end{align}

\begin{align}\label{e.L_0=}
    L_0 : = \lim_{\beta\to\infty}\int_0^1\xi''(t)\gamma_{P,\beta}(t)\d t
\end{align}

\begin{align}\label{e.nu_0=}
    \d \nu_0(t):=\mathds{1}_{[0,1)}(t)\gamma_0(t)\d t + \frac{1}{\xi''(1)}\Ll(L_0 - \int_0^1\xi''(t)\gamma_0(t)\d t\Rr)\d\boldsymbol{\delta}_1( t)
\end{align}

\begin{align}\label{e.int..phi->int..phidnu}
    \lim_{\beta\to\infty}\int_0^1\xi''(t)\gamma_{P,\beta}(t)\phi(t)\d t= \int_0^1\xi''(t)\phi(t)\d \nu_0(t).
\end{align}

\begin{align}
    \Psi_{\gamma,\beta}(0,x) = \sup_{\alpha\in\prog_1}\bigg(\frac{1}{\beta}\E\log\cosh \beta\Ll( x + \int_0^1\xi''(t)\gamma(t)\alpha_t\d t +\int_0^1\sqrt{\xi''(t)}\d W_t\Rr)\notag
    \\
    -\frac{1}{2}\int_0^1 \xi''(t)\gamma(t)\E[\alpha^2_t]\d t\bigg) \label{e.AC_rep_1}
\end{align}

\begin{align}\label{e.AC_rep_2}
\begin{split}
    \Psi_{\gamma}(0,x) = \sup_{\alpha\in\prog_1}\bigg(\E\Ll| x + \int_0^1\xi''(t)\gamma(t)\alpha_t\d t +\int_0^1\sqrt{\xi''(t)}\d W_t\Rr|
    \\
    -\frac{1}{2}\int_0^1 \xi''(t)\gamma(t)\E[\alpha^2_t]\d t\bigg)
\end{split}
\end{align}

\begin{lemma}[Lower bound]\label{l.laplace_>}
For every $s \in (0,1)$, we have
\begin{equation}
\label{e.laplace_>}
\lim_{N\to \infty} \frac{1}{sN} \log \E \Ll[ e^{sN L_N} \Rr]   \geq  \inf_{\gamma \in \cM_{s }}\bigg( \frac 1 s \log \E\Ll[e^{s\Psi_{\gamma}(0,g_1+h_1)}\Rr] - \frac{1}{2} \int_0^1 t \xi''(t) \gamma(t) \, \d t \bigg).
\end{equation}
\end{lemma}

\begin{proof}
For each $n\in \N$, let $\sign_n:\R\to[-1,1]$ be defined by
\begin{align*}
    \sign_n(x) = 
    \begin{cases}
        1, \qquad &\text{if }x\geq 0,
        \\
        2nx + 1, \qquad &\text{if } -\frac{1}{n}\leq x<0,
        \\
        -1, \qquad &\text{if } x<-\frac{1}{n}.
    \end{cases}
\end{align*}
We view $(\sign_n)_{n\in\N}$ as approximations of the sign function, namely,
\begin{align*}
    \lim_{n\to\infty} \sign_n(x) = \sign(x):=
    \begin{cases}
        1  , \qquad &\text{if }x\geq 0,
        \\
        -1 , \qquad &\text{if }x< 0.
    \end{cases}
\end{align*}
Let $\gamma_0$ be given as above~\eqref{e.L_0=}. We write $\bar h= g_1+h_1$ for brevity and recall that $\bar h$ is independent from the Wiener process $W$.
Fix any $\alpha\in\prog_1$. For every $\tau\in(0,1)$ and $n\in\N$, define
\begin{align*}
    \phi_{\tau,n}(t) = \alpha_t\mathds{1}_{[0,\tau)}(t) + \sign_n\Ll(\bar h+\int_0^t\xi''(r)\gamma_0(r)\alpha_r\d r + \int_0^t\sqrt{\xi''(r)}\d W_r\Rr)\mathds{1}_{[\tau,1]}(t)
\end{align*}
for every $t\in[0,1]$. Then, we have $\phi_{\tau,n}\in\prog_1$ and $\lim_{t\to1-}\phi_{\tau,n}(t) = \phi_{\tau,n}(1)$, since $g_n$ is continuous. We also have
\begin{align}\label{e.phi_n...}
\begin{split}
    \phi_n(t) &:=\lim_{\tau\to1-}\phi_{\tau,n}(t)=\alpha_t\mathds{1}_{[0,1)}(t)+\sign_n(X)\mathds{1}_{\{1\}}(t),
    \\
    \lim_{n\to\infty}\phi_n(t) &= \alpha_t\mathds{1}_{[0,1)}(t)+\sign(X)\mathds{1}_{\{1\}}(t)
\end{split}
\end{align}
for every $t\in[0,1]$, where we have set
\begin{align}\label{e.X=}
    X:=\bar h +\int_0^1\xi''(t)\gamma_0(t)\alpha_t\d t+ \int_0^1\sqrt{\xi''(t)}\d W_t.
\end{align}
Recall that $\gamma_{P,\beta}$ is the minimizer of~\eqref{e.lim1/sNlogE[Z_N(beta)^s/beta]}. We write $\spadesuit= \lim_{N\to \infty} \frac{1}{sN} \log \E \Ll[ e^{sN L_N} \Rr] $ for brevity. Then, we have
\begin{align*}
    \spadesuit\stackrel{\eqref{e.lim1/sNloge^sNL_N}\eqref{e.lim1/sNlogE[Z_N(beta)^s/beta]}}{=}\lim_{\beta\to\infty} \bigg( \frac{1}{s}\log\E\Ll[e^{\Psi_{\gamma_{P,\beta},\beta}(0,\bar h)}\Rr] - \frac{1}{2} \int_0^1 t \xi''(t) \gamma_{P,\beta}(t) \, \d t \bigg)
    \\
    \stackrel{\eqref{e.AC_rep_1}}{\geq}\lim_{\beta\to\infty}\bigg(\frac{1}{s}\log\E_{\bar h}\Ll[e^{\frac{s}{\beta}\E_{W}\log\cosh \beta\Ll( \bar h + \int_0^1\xi''(t)\gamma_{P,\beta}(t)\phi_{\tau,n}(t)\d t +\int_0^1\sqrt{\xi''(t)}\d W_t\Rr)}\Rr]
    \\
    -\frac{1}{2}\int_0^1 \xi''(t)\gamma_{P,\beta}(t)\Ll(\E[\phi_{\tau,n}^2(t)]+t\Rr)\d t\bigg)
    \\
    \stackrel{\text{(Fatou's lemma)\eqref{e.int..phi->int..phidnu}}}{\geq} \frac{1}{s}\log\E_{\bar h}\Ll[e^{s\E_W\Ll| \bar h + \int_0^1\xi''(t)\phi_{\tau,n}(t)\d\nu_0(t) +\int_0^1\sqrt{\xi''(t)}\d W_t\Rr|}\Rr]
    \\
    -\frac{1}{2}\int_0^1 \xi''(t)\Ll(\E[\phi_{\tau,n}^2(t)]+t\Rr)\d \nu_0(t)
\end{align*}
where $\E_W$ averages over $\phi_{\tau,n}$ and $W$ and $\E_{\bar h}$ averages over $\bar h$.
Using~\eqref{e.nu_0=}, \eqref{e.phi_n...}, and~\eqref{e.X=} together with the dominated convergence theorem, we send $\tau\to1-$ and then $n\to\infty$ to get
\begin{align*}
    \spadesuit & \geq \frac{1}{s}\log\E_{\bar h}\Ll[e^{s\E_W\Ll|X + \sign(X)\xi''(1)\nu_0(1)\Rr|}\Rr]
    \\
    &-\frac{1}{2}\int_0^1 \xi''(t)\Ll(\E[\alpha^2_t]+t\Rr)\gamma_0(t)\d t - \frac{1}{2}\Ll(\E\Ll[\sign(X)^2\Rr]+1\Rr)\xi''(1)\nu_0(1).
\end{align*}
Since
\begin{align*}
    &\Ll|X+\sign(X)\xi''(1)\nu_0(1)\Rr|
    \\
    &=\Ll(X+\xi''(1)\nu_0(1)\Rr)\mathds{1}_{X>0}- \Ll(X-\xi''(1)\nu_0(1)\Rr)\mathds{1}_{X<0}+\xi''(1)\nu_0(1)\mathds{1}_{X=0}
    \\
    & =|X|+\xi''(1)\nu_0(1)
\end{align*}
and $\Ll(\E\Ll[\sign(X)^2\Rr]+1\Rr)\xi''(1)\nu_0(1)=2\xi''(1)\nu_0(1)$, we obtain
\begin{align*}
    \spadesuit \geq \frac{1}{s}\log\E_{\bar h}\Ll[e^{s\E_W\Ll|X\Rr|}\Rr]
    -\frac{1}{2}\int_0^1 \xi''(t)\Ll(\E[\alpha^2_t]+t\Rr)\gamma_0(t)\d t .
\end{align*}
Recall the expression of $X$ in~\eqref{e.X=} and that $\alpha\in\prog_1$ is chosen arbitrarily. Using these combined with~\eqref{e.AC_rep_2}, we get
\begin{align*}
    \spadesuit \geq \frac{1}{s}\log\E_{\bar h}\Ll[e^{s\Psi_{\gamma_0}(0,\bar h)}\Rr]
    -\frac{1}{2}\int_0^1 t\xi''(t)\gamma_0(t)\d t
\end{align*}
which implies~\eqref{e.laplace_>}.
\end{proof}

\end{document}
